\documentclass[a4paper,12pt,oneside]{book} % Formato documento: foglio a4, font size 12, solo fronte
\usepackage[utf8]{inputenc} % Codifica caratteri
\usepackage[lmargin=2cm,rmargin=2cm,tmargin=2.5cm,bmargin=2cm]{geometry} % Bordi pagina come da formato tesi Unibs
\linespread{1.5} % Interlinea
\setlength{\parindent}{0pt} % No spazio prima dei paragrafi
\setlength{\parskip}{1em}

\usepackage[backend=biber, sorting=none]{biblatex} % Per poter scrivere la bibliografia su un file a parte
\addbibresource{bibliography.bib} % File contentente la bibliografia

\usepackage{titlesec} % Per modificare il formato dei capitoli
\titleformat{\chapter}[hang]{\normalfont\Huge\bfseries}{\thechapter.}{8pt}{\Huge\bfseries}
\titlespacing*{\chapter}{0pt}{0pt}{20pt}

\usepackage{fancyhdr} % Per modificare l'intestazione delle pagine
\pagestyle{fancy} % Customizzazione testata e piè di pagina
\renewcommand{\chaptermark}[1]{\markboth{#1}{#1}}
\fancyhf{}
\fancyhead[L]{\thechapter. \leftmark}
\fancyfoot[C]{\thepage}

\usepackage{graphicx} % Per poter inserire immagini
\usepackage{float} % Per controllare il posizionamento di immagini e tabelle
\usepackage{amsmath}
\usepackage{hyperref} % Per scrivere equazioni

\title{Implementation of The Multiple Knapsack Problem in the Kernel Search}
\author{Matteo Beatrice, Luca Cotti, Giacomo Golino}
\date{December 2020}

\begin{document}
    %Frontespizio
    \begin{titlepage}
    \begin{figure}[H] % Logo Unibs
        \centering
        \includegraphics[width=72.4mm]{images/logo_unibs}
    \end{figure}

    \begin{center}
        \LARGE{\uppercase{Università degli Studi di Brescia}}\\ % Doppio backslash = vai a capo
        \vspace{5mm} % Spazio verticale
        \large{\uppercase{Dipartimento di Ingegneria dell'informazione}}\\
        \vspace{5mm}
        \large{Corso di Laurea Magistrale in Ingegneria Informatica}\\
    \end{center}

    \vspace{10mm}

    \begin{center}
        \LARGE{\textbf{Implementation of The Multiple Knapsack Problem in the Kernel Search}}\\
    \end{center}

    \vspace{10mm}

    \begin{flushleft}
        \large
        \textbf{Docente:}\\
        Ch.ma Prof.sa Renata Mansini
    \end{flushleft}

    \begin{flushright}
        \large
        \textbf{Studenti:}\\
        Matteo Beatrice (739848)\\
        Luca Cotti (719204)\\
        Giacomo Golino (719210)
    \end{flushright}

    \vspace*{\fill} % Spazio bianco per mettere il testo in fondo alla pagina

    \rule{0.8\textwidth}{0.6pt}\\ % Linea nera
    \centering{\Large{Anno Accademico 2021/2022}}
\end{titlepage}

    %Corpo
    \tableofcontents
    \hyphenpenalty=10000 % Non spezzare le parole a capo

    \chapter{Introduction}
    \chapter{The Multiple Knapsack Problem}
Given a set of \(m\) knapsacks with a known \(\text{capacity } c_{i}\) \((i=1,\dots,m)\) and a set of \(n\) items with
known \(\text{profit } p_{j}\) and \(\text{weight } w_{j}\) \((j=1,\dots,n)\) the Multiple Knapsack Problem (MKP) consists in
selecting \(m\) disjoint subsets of items (one subset per knapsack) such that the total weight of the items in the knapsack
does not exceed its capacity, and the profit of the selected items is maximized.

\section{Mathematical model}
There are quite a few formulations for the MKP, with different degrees of
complexity and performance.
The formulation that was used in this project is the classical and most intuitive
one, that uses the following binary variables:
\[x_{ij}=
\begin{cases}
    1 & \quad \text{if item } j \text{ is packed into knapsack } i\\
    0 & \quad \text{otherwise}
\end{cases}
\]

The model is as follows:
\begin{equation}
    \label{eq:obj}
    \max{z} = \sum_{i=1}^{n} \sum_{j=1}^{n} p_{j} x_{ij}
\end{equation}
\begin{equation}
    \label{eq:cap}
    \sum_{j=1}^{n} w_{j} x_{ij} \leq c_{i} \qquad i=1,\dots,m
\end{equation}
\begin{equation}
    \label{eq:sel}
    \sum_{i=1}^{m} x_{ij} \leq 1 \qquad j=1,\dots,n
\end{equation}
\begin{equation}
    x_{ij} \in \{0,1\} \qquad i=1,\dots,m \quad j=1,\dots,n
\end{equation}

The objective function~\eqref{eq:obj} maximises the profit of the items in the knapsacks.\\
Constraints~\eqref{eq:cap} impose that the capacity of each knapsack is respected, while constraints~\eqref{eq:sel}
ensure that each item is packed in no more than one knapsack.

\section{Real life examples}
The MKP has various real life examples, in various fields.

The most straightforward example might be the cargo organization for a transport company.
In this scenario, the knapsack may be the trucks or containers.
Their capacity is the effective capacity of the truck/container.
The items are the objects that must be delivered.
Their weight is the actual weight of the item, while the profit is the priority of the delivery.

Another example might be the project selection:
The knapsacks may be the years of development of the project, and their capacity is the available budget for
that year.
The items can be the available projects: the weight is the cost of the project, the profit is the monetary gains
from the project completion.
    \chapter{Implementation and Test Instances}

\section{Kernel Search}
The code for the kernel search was provided by the teacher.\\
The code implements a generic kernel search, with the following default paramenters:
\begin{itemize}
    \item Presolve = 2
    \item Time limit = 30 seconds
    \item Threads = 12
    \item Sort variables by profit and reduced cost
    \item Kernel size = 15\% of the number of variables
    \item Bucket size = 2.5\% of the number of variables
    \item Number of iterations = 2
    \item Time limit for each bucket = 5
\end{itemize}

The code was modified to implement the MKP model, and it was then
refactored to improve its quality and efficiency.
The kernel search algorithm itself and the default parameters were not changed
at this time.
The updated code can be found at the GitHub repository
\url{https://github.com/Golino98/KernelSearchGolinoCottiBeatrice}.

\section{Exact Solver}
\section{Test Instances}
The
    \printbibliography
\end{document}